\section{Introducción}

Este trabajo\footnote[1]{Se utilizan los Lineamientos FACSO para la elaboración de este proyecto, donde los proyectos "(...) pueden seguir un formato de investigación empírica, discusión teórica, o intervención político social. El desarrollo del Seminario de Grado va a depender del tipo de enfoque que la o el estudiante decida adoptar. No obstante, hay elementos comunes que deberán ser abordados a lo largo de las secciones del Seminario de Grado. Estos incluyen antecedentes problematización, objetivos, perspectiva teórica y metodológica. El orden en que se presentan estos elementos puede variar (por ejemplo, la problematización puede situarse después de los antecedentes y perspectiva teórica)."} analiza comparativamente la relación entre los ingresos monetarios de los hogares y el costo de una canasta básica de referencia en Chile, Brasil y Argentina. La posibilidad de cubrir necesidades esenciales depende tanto de los ingresos y la inserción laboral como de la evolución de los precios, la composición familiar y las instituciones de protección social \parencite{cepalPanoramaSocialAmerica2025, GindlingRonconi2025, DuranKremerman2024}. La inflación y el encarecimiento de bienes básicos afectan de manera desigual a los hogares según su posición en la distribución del ingreso y sus estrategias de consumo \parencite{Adam2025, Nessieretal2026}.

El análisis utiliza las fuentes disponibles en cada país: CASEN 2024 y la IX Encuesta de Presupuestos Familiares (EPF 2021--2022) para Chile; PNAD Contínua 2025 y primer trimestre de 2026 para Brasil; y EPH del primer trimestre de 2026 para Argentina. Estos períodos no son simultáneos ni conforman una serie histórica común. La comparación se plantea como una descripción por país, encuesta y cobertura territorial, preservando las diferencias de definición de ingreso, diseño muestral y medición del costo de vida. El flujo inicial utiliza CASEN con cobertura nacional, restringe PNAD 2025 a capitales estatales y EPH 1T2026 al Gran Buenos Aires; el archivo PNAD 1T2026 aún requiere su diccionario específico. La EPF chilena permite describir gasto observado, pero no reemplaza el valor de una canasta básica oficial. En Brasil, el DIEESE publica una canasta mensual por capital que constituye una referencia independiente, no una estadística gubernamental; en Argentina, la CBA/CBT del INDEC tiene como referencia el Gran Buenos Aires y se expresa por adulto equivalente.

La pregunta central es cuánto del costo de una canasta básica local puede cubrir el ingreso monetario de los hogares y cómo varía esa capacidad según sus características sociales. La razón entre ingreso y costo solo se calculará cuando ambas medidas correspondan a unidades, períodos y coberturas territoriales compatibles. Este indicador permite describir capacidad de cobertura dentro de cada caso, pero no convierte las canastas nacionales en bienes equivalentes ni autoriza una clasificación simple de bienestar entre países.

La investigación articula antecedentes sobre distribución del ingreso \parencite{ResumenResultadosCASEN2023}, protección social \parencite{Stampinietal2025, Carrere2026}, instituciones laborales \parencite{BarreraInsuaNoguera2026, DuranS2026} y reproducción material de los hogares. Su marco teórico integra la economía política del trabajo y la acumulación de capital \parencite{DuranSilva2022, MarxElCapital2008} con enfoques sobre estratificación social y bienestar en América Latina \parencite{cepalPanoramaSocialAmerica2025, CostaRibeiroCarcalhaes2024}.